\usepackage{fontspec}
\usepackage[english, russian]{babel}
\usepackage{graphicx}
\usepackage[a4paper,
            left=3cm,
            right=3cm,
            top=2cm,
            bottom=3cm]{geometry}
\usepackage{amsmath, amsfonts}
\usepackage{tabularray}

\usepackage{pgfplots}
\usepackage{tikz}
  \usetikzlibrary{positioning, shapes.geometric, calc, arrows.meta, chains, shapes.symbols, shapes.misc, patterns}

\setmainfont{CMU Serif}
\setsansfont{CMU Sans Serif}
\setmonofont{CMU Typewriter Text}

\usepackage{hyperref}
\hypersetup{
    colorlinks=true,        % цветные ссылки (false - черные рамки)
    linkcolor=black,        % цвет ссылок на разделы
    urlcolor=blue,          % цвет URL-ссылок
    citecolor=blue,         % цвет ссылок на литературу
    pdfborder={0 0 0}       % убрать рамки вокруг ссылок
}
\addto\extrasrussian{%
  \def\figureautorefname{рис.}%
  \def\tableautorefname{табл.}%
  \def\sectionautorefname{разд.}%
  \def\equationautorefname{ур-ние}%
}

\usepackage{float}
\newfloat{scheme}{htbp}{los}
\floatname{scheme}{Схема}

\usepackage[justification=centering]{caption}
\DeclareCaptionLabelSeparator{emdash}{\space\textemdash\space}
\captionsetup{labelsep=emdash}

\usepackage{listings}
% Настройки для листингов кода
\lstset{
    language=C++,
    basicstyle=\ttfamily\small,
    keywordstyle=\color{blue},
    commentstyle=\color{green!60!black},
    stringstyle=\color{red},
    numbers=left,
    numberstyle=\tiny,
    stepnumber=1,
    numbersep=5pt,
    frame=single,
    tabsize=4,
    breaklines=true,
    showstringspaces=false,
    captionpos=b
}

\input{sections/Block}

\newcommand{\abs}[1]{\left|#1\right|}

\usepackage[
    backend=biber,
    style=gost-numeric,
    sorting=none
]{biblatex}
\addbibresource{refs.bib}